\documentclass{../Templates/sbs-exam}

\usepackage{ulem}
\usepackage{array}
\usepackage{enumitem}
\usepackage{tabularx}
\usepackage{ulem}
\usepackage{xcolor}
\usepackage{multirow}
\usepackage{multicol}
\usepackage{graphicx} % 图形支持包（必须）
\usepackage{fancyhdr} % 页眉控制包（必须）
\usepackage{titling} % 提供高级标题控制
\usepackage{lmodern} % 现代字体支持
\usepackage{amsmath, amssymb} % 数学符号和公式支持
\usepackage{lipsum} % 用于生成示例文本
\usepackage{caption} % 图表标题支持

% ===== 资源路径配置 =====
\graphicspath{
  {../Assets/}     % 公共资源
  {../Media/}      % 试卷专用资源
  {./Figures/} % 本试卷资源
}

% ===== 考试元数据配置 =====
\examtitle{Monthly Test, 2025-2026 T2}
\examsubject{G10 A3 Mathematics}{Mar 2026}{60 Minutes}{Shi Feng}
\MarksBreakdown{2 \text{ Marks}}{10}{2 \text{ Marks}}{20}{30 \text{ Marks}}{90 \text{ Marks}}

% 自定义命令用于处理记分内容
\newcommand{\imarks}[1]{\hfill \textbf{(#1 marks)}}
\newcommand{\totalmarks}[1]{\par\noindent\hfill \textbf{(Total for question = #1 marks)}\par}

\begin{document}

\maketitle

\examnotice

\makemarksheet

\examtools

\newpage

\makeatletter

\section*{Multiple Choice Questions (10 questions, 2 marks each)}
\begin{enumerate}
    \item The midpoint of the line segment joining $A(-3,4)$ and $B(5,2)$ is
    \begin{enumerate}
        \item $(1,3)$
        \item $(4,1)$
        \item $(2,2)$
        \item $(1,2)$
    \end{enumerate}
    
    \item The equation of the circle with centre $(0,0)$ and radius $5$ is
    \begin{enumerate}
        \item $x^2 + y^2 = 25$
        \item $x^2 + y^2 = 5$
        \item $(x-5)^2 + (y-5)^2 = 0$
        \item $x^2 + y^2 = 10$
    \end{enumerate}
    
    \item $\log_2 16$ is equal to
    \begin{enumerate}
        \item $2$
        \item $4$
        \item $8$
        \item $16$
    \end{enumerate}
    
    \item Which of the following is a geometric sequence?
    \begin{enumerate}
        \item $2,5,8,11,\dots$
        \item $3,6,12,24,\dots$
        \item $1,4,9,16,\dots$
        \item $10,7,4,1,\dots$
    \end{enumerate}
    
    \item The value of $\sin 150^\circ$ is
    \begin{enumerate}
        \item $\frac{1}{2}$
        \item $-\frac{1}{2}$
        \item $\frac{\sqrt{3}}{2}$
        \item $-\frac{\sqrt{3}}{2}$
    \end{enumerate}
    
    \item In the expansion of $(1+x)^5$, the coefficient of $x^2$ is
    \begin{enumerate}
        \item $5$
        \item $10$
        \item $1$
        \item $2$
    \end{enumerate}
    
    \item The sum of the first $10$ terms of the arithmetic series $3+7+11+15+\dots$ is
    \begin{enumerate}
        \item $210$
        \item $220$
        \item $230$
        \item $240$
    \end{enumerate}
    
    \item Which of the following is equivalent to $\tan\theta$?
    \begin{enumerate}
        \item $\frac{\sin\theta}{\cos\theta}$
        \item $\frac{\cos\theta}{\sin\theta}$
        \item $\frac{1}{\sin\theta}$
        \item $\frac{1}{\cos\theta}$
    \end{enumerate}
    
    \item The circle $x^2 + y^2 - 4x + 6y - 12 = 0$ has centre
    \begin{enumerate}
        \item $(-2,3)$
        \item $(2,-3)$
        \item $(-2,-3)$
        \item $(2,3)$
    \end{enumerate}
    
    \item The value of $\left(\frac{1}{2}\right)^{-2}$ is
    \begin{enumerate}
        \item $\frac{1}{4}$
        \item $-\frac{1}{4}$
        \item $4$
        \item $-4$
    \end{enumerate}
\end{enumerate}

\section*{Short Answer Questions (20 questions, 2 marks each)}

\subsection*{True/False}
State whether each statement is True or False.

\section*{True or False Questions (10 questions)}
\begin{enumerate}
    \item The perpendicular bisector of a chord always passes through the centre of the circle. \hfill True \quad False
    
    \item $\log_a 1 = 0$ for any $a>0$, $a\neq 1$. \hfill True \quad False
    
    \item In a geometric sequence with common ratio $r$, if $|r|<1$, the sum to infinity exists. \hfill True \quad False
    
    \item The graph of $y = a^x$ (with $a>1$) passes through the point $(0,0)$. \hfill True \quad False
    
    \item $\sin^2\theta + \cos^2\theta = 1$ for all $\theta$. \hfill True \quad False
    
    \item The sum of the first $n$ natural numbers is $\frac{n(n+1)}{2}$. \hfill True \quad False
    
    \item The coefficient of $x^3$ in the expansion of $(1-2x)^5$ is $-80$. \hfill True \quad False
    
    \item The tangent to a circle at a point is perpendicular to the radius at that point. \hfill True \quad False
    
    \item $\log_2 8 = 3$ and $\log_2 \frac{1}{8} = -3$. \hfill True \quad False
    
    \item The equation $(x-3)^2 + (y+2)^2 = 9$ represents a circle with centre $(-3,2)$ and radius $9$. \hfill True \quad False
\end{enumerate}

\section*{Fill in the Blanks (10 questions)}
\begin{enumerate}
    \item The midpoint of $A(2,-5)$ and $B(6,1)$ is \underline{\hspace{2cm}}.
    
    \item The equation of the circle with centre $(4,-3)$ and radius $2$ is \underline{\hspace{2cm}}.
    
    \item The value of $\log_3 27$ is \underline{\hspace{2cm}}.
    
    \item The common ratio of the geometric sequence $4, -8, 16, -32, \dots$ is \underline{\hspace{2cm}}.
    
    \item The exact value of $\cos 45^\circ$ is \underline{\hspace{2cm}}.
    
    \item The sum to infinity of the geometric series $1 + \frac{1}{2} + \frac{1}{4} + \frac{1}{8} + \dots$ is \underline{\hspace{2cm}}.
    
    \item The expansion of $(2+x)^3$ is \underline{\hspace{2cm}}.
    
    \item The gradient of the line joining $(1,2)$ and $(5,10)$ is \underline{\hspace{2cm}}.
    
    \item In the unit circle, $\sin\theta$ is equal to the \underline{\hspace{2cm}}-coordinate of the point on the circle.
    
    \item The solution to $2^{x} = 32$ is $x =$ \underline{\hspace{2cm}}.
\end{enumerate}

\newpage

\section*{Long Answer Questions}

\begin{enumerate}

\item % Q1: Coordinate geometry (circle, tangent, chord)
The circle $C$ has equation $(x-3)^2 + (y+2)^2 = 25$.
\begin{enumerate}
    \item Find the centre and radius of $C$. \imarks{4}
    \item Show that the point $P(6,2)$ lies on $C$. \imarks{2}
    \item Find the equation of the tangent to $C$ at $P$. Give your answer in the form $ax+by+c=0$, where $a,b,c$ are integers. \imarks{4}
\end{enumerate}

\vspace{8cm}

\item % Q2: Exponentials and logarithms (solving equations)
\begin{enumerate}
    \item Solve $3^{2x-1} = 27$. \imarks{3}
    \item Solve $\log_2(x+3) - \log_2(x-1) = 2$. \imarks{4}
    \item Given that $\log_a 4 = 0.5$, find the value of $a$. \imarks{3}
\end{enumerate}

\vspace{12cm}

\item % Q3: Sequences and series (arithmetic/geometric, sum to infinity)
A geometric series has first term $a$ and common ratio $r$, where $0<r<1$. The sum of the first two terms is $16$ and the sum to infinity is $18$.
\begin{enumerate}
    \item Show that $r = \frac{1}{3}$ and find $a$. \imarks{4}
    \item Find the sum of the first $5$ terms of the series, giving your answer as a fraction. \imarks{3}
    \item Find the least value of $n$ such that $S_n > 17.9$. \imarks{3}
\end{enumerate}


\end{enumerate}

\end{document}